\clearpage
\section{Related Work}
\textbf{What to update.} RialTo reconstructs a scene, transfers demonstrations into the twin, and uses simulation reinforcement learning to improve a real-world imitation policy \cite{rialto}. MimicGen adapts demonstrations to new configurations \cite{mandlekar2023mimicgen}; ARSTAG uses agents to generate task-specific robot data \cite{arstag}. Agentic RSR already uses execution feedback in scene construction and policy development \cite{agenticrsr}. Our distinction is not the use of feedback itself. We hold the packing task fixed and compare spending a post-deployment round on the scene model against spending it on demonstrations or policy training.

\textbf{What to label.} DAgger changes the state distribution on which an expert supplies actions; HG-DAgger lets a human decide when to intervene \cite{ross2011dagger,kelly2019hgdagger}. LoopWAM is especially close: it collects deployment trajectories, annotates quality and retrains with expert and historical rollout data. On RoboChallenge Table30-V2, success rose from 23.67\% to 43.33\% over three rounds \cite{loopwam}. That result establishes the value of learning from mixed-quality physical rollouts. It does not compare collecting another batch with spending that round correcting geometry. Our deployment-label arm addresses that comparison in simulation; it is not a reproduction of LoopWAM.

\textbf{What to measure.} Domain randomization varies the training distribution \cite{tobin2017dr}. SimOpt adapts that distribution using real experience \cite{chebotar2019asid}, and ASID selects informative interactions for system identification \cite{asid}. These approaches motivate alternatives to assuming a known model. REFLECT summarizes executions to explain failures and guide replanning \cite{reflect}; AHA uses failure explanations to refine rewards, plans and verification \cite{aha}; CycleVLA backtracks at the subtask level \cite{cyclevla}. Rather than evaluate the explanation or the next motion, we measure the return on the next round of model or policy updates.

\section{Decision and Evaluation Protocol}
Let $W$ be a deployment scene and $T$ its reconstruction. A policy $\pi$ acts using features computed with $T$; success is scored in $W$. Both scenes use MuJoCo \cite{mujoco}. A bank is a set of task instances, not a robot. The task is to pack four objects in sequence using a hierarchical controller and a 55-dimensional state representation. The selector receives no cause identifier. Its inputs include observed support-height residuals, the distribution of failed subgoals, recent success and the history of updates. The injected cause is available only for analysis and to a reference selector.

At round $t$, a selector chooses an operation $o_t$ from this information. We compare final complete-task success rates (SR) under limits on feedback episodes, training transitions and rounds, recording actual expenditure rather than equating a label with a deployment. Standard initialization uses 80,000 labeled transitions. Each of three rounds provides 24 feedback episodes; a data operation adds up to 40,000 labels and retraining uses 8,000 optimization steps. Calibration uses observations but adds no training labels. Stopping saves subsequent rounds. The open-loop control spends its 200,000-label allowance at initialization. Feedback episodes are not themselves 40,000 expert labels, and deployment-side expert data would incur a different physical acquisition cost.

The geometric condition sets believed support height to 0.46 of the true height (9.2 versus 20 mm). Under limited coverage, training layouts occupy half the workspace and deployment uses the full layout distribution. The recovery condition omits recovery demonstrations and tests under perturbation; the healthy condition introduces neither a reconstruction error nor a training deficit. Healthy thus means no injected defect, not guaranteed success. Only recovery uses perturbed evaluation in the four-condition comparison. Differences between methods within a condition have a cleaner interpretation than differences between conditions.
